\section{Hyperbolicity from two successive fillings}\label{sec:geometry}

\subsection{A circle around the moving strands}

Set \(\rho_b=\sigma_{b-1}\cdots\sigma_4\) and
\(\theta_d=\sigma_2\cdots\sigma_{b-2}\). The braid relation gives
\(\rho_b\delta_b\rho_b^{-1}=\theta_d\). For instance,
conjugation of \(\sigma_3\) by \(\sigma_4^{-1}\) gives
\(\sigma_3\sigma_4\sigma_3^{-1}\) when \(b=5\).
In the step from \(b\) to \(b+1\), the new letter
\(\sigma_{b-1}\) of \(\theta_{d+1}\) is conjugated by
\(\rho_b^{-1}\sigma_b^{-1}\), giving \(\sigma_b\) because
\(\sigma_b^{-1}\sigma_{b-1}\sigma_b
=\sigma_{b-1}\sigma_b\sigma_{b-1}^{-1}\).
The earlier letters are unchanged by \(\sigma_b\) before conjugating
by \(\rho_b^{-1}\). It follows that the two braid words
\[
 Q_b=\rho_bU\rho_b^{-1}\theta_d^{d+1},
 \qquad Q'_b=\rho_bU\rho_b^{-1}\theta_d
\]
are convenient starting points for a twist on the \(d\) consecutive
strands in positions \(2,\ldots,b-1\).

Add strand \(b+1\) and append
\[
       P_b=\sigma_b\sigma_{b-1}\cdots\sigma_3\sigma_2^2
                         \sigma_3\cdots\sigma_{b-1}\sigma_b^{-1}.
\]
The last strand closes to an unknot \(C\), bounding a disk pierced
once by each strand in positions \(2,\ldots,b-1\) and missing the
other strands. One way to see the disk is to factor \(P_b\) as
\(A_{2,b+1}\cdots A_{b-1,b+1}\), where
\[
  A_{i,b+1}=\sigma_b\cdots\sigma_{i+1}\sigma_i^2
                          \sigma_{i+1}^{-1}\cdots\sigma_b^{-1}.
\]
Each factor pushes the last puncture once around puncture \(i\).
The stems in this ordered product run along the same side of the
adjacent punctures, so their concatenation bounds an oval around
precisely those punctures. Equivalently, in the closure of the last
strand, its two opposite crossings across strand \(b\) make a
cancelling detour. The remaining oval passes over the selected strands
on the outward trip and under them on the return. This gives the claimed
disk in the slab of \(P_b\), including the two outside positions.
The knot component of \(\cl{Q_bP_b}\) has the permutation already
computed in Section~\ref{sec:color}; the same is true for \(Q'_b\).

For a disk-bounding circle with this marking, filling along
\((1,-q)\) means killing \(\mu_C-q\lambda_C\), where \(\mu_C\)
is a positive meridian and \(\lambda_C\) the preferred disk
longitude. A disk twist on the solid torus obtained by removing \(C\)
takes this slope to \(\mu_C\). Localizing the twist to the slab
inserts \(q\) full twists on the disk-piercing strands. To check the
sign, put \(w=a_1\cdots a_d\) for the selected meridians and \(c\)
for the circle meridian. Tracing the outward and return crossings
of the oval gives the selected outputs \(c^{-1}a_i c\) and circle
output \((wc)^{-1}c(wc)\). The return has crossed the whole ordered
tuple, which contributes \(w\). The two outside crossings conjugate
the outside meridian by inverse powers of the initial and final
circle labels and cancel after closure. The closure relation is
\([c,w]=1\), and \(\lambda_C\) is represented by \(w\), the
boundary of the punctured disk. Filling gives \(c=w^q\).
By the Hurwitz formula, \(\theta_d^d\) conjugates the whole
selected tuple by \(w^{-1}\), so the geometric disk twist inserts
\(\theta_d^{dq}\) in our convention.

Let \(M_b=\Ext(\cl{Q_bP_b})\) and \(M'_b=\Ext(\cl{Q'_bP_b})\),
where \(\Ext(L)\) denotes the link exterior. The preceding disk
twist identifies \(M_b\) with \(M'_b\) as exteriors without a
fixed peripheral basis, since their knot braids differ by one full
twist on the disk. In the original marking of \(M_b\) we have
\begin{equation}\label{eq:kfilling}
      M_b\bigl(C:(1,-d(d+1)k)\bigr)\cong \Ext(K_{b,k}).
\end{equation}
Indeed, the inserted power of \(\theta_d\) is \(d^2(d+1)k\).
Conjugating the resulting braid by \(\rho_b^{-1}\) gives the word
defining \(K_{b,k}\).

\subsection{A single exterior for the first filling}

Consider the six-strand braid
\[
 R=\sigma_4\sigma_3\sigma_2^2\sigma_3\sigma_4\sigma_5^{-1}
                 \sigma_4 U\sigma_4^{-1}\sigma_2\sigma_3\sigma_5
\]
and the seven-strand suffix \(D=\sigma_6\sigma_5\sigma_4^2
\sigma_5\sigma_6^{-1}\). Set \(N=\Ext(\cl{RD})\).
The closure of \(R\) has two components, a knot component containing
\(U\) and a circle. The fourth and fifth closure arcs lie on these
different components. These assignments also follow by tracing the
strand permutation of \(R\), which fixes position five and cycles
the other five positions. The new last component \(D_0\) bounds a
disk around positions four and five, by the same oval
construction as for \(P_b\). Thus \(N\) has three boundary tori,
and filling the new circle along the disk-marked slope gives
\begin{equation}\label{eq:nfilling}
      N\bigl(D_0:(1,-q)\bigr)\cong\Ext(\cl{R\sigma_4^{2q}})
                         \qquad(q\in\mathbb Z).
\end{equation}

\begin{lemma}\label{lem:isotopy}
For every \(b\ge5\), the oriented link \(\cl{Q'_bP_b}\) is isotopic to
\(\cl{R\sigma_4^{2(b-5)}}\), with the knot and circle components
identified.
\end{lemma}

Here is how the isotopy works. After braid conjugation and
commutation,
\[
  \sigma_3^{-1}(U\delta_b)\sigma_3
    =(\sigma_3^{-1}U\sigma_2\sigma_3)\sigma_4\sigma_5\cdots\sigma_{b-1}.
\]
On the knot component alone this is an iterated positive
stabilization. In the link, the disk of \(C\) meets the additional
strands, so the corresponding curls cannot simply be removed.
First cancel the detour of the circle and the two sides of the
conjugating braid. The outer crossings then form a ladder along the
circle. Slide each successive triangle through the already formed
clasps by two Reidemeister III moves per clasp. At its end a
Reidemeister III move exposes a curl, and a Reidemeister I move
leaves one more two-crossing positive clasp. Repeating this leaves
the \(b=5\) diagram with \(b-5\) clasps in the band of the fourth and
fifth closure arcs of \(R\). Appendix~\ref{app:ports} proves the
isotopy by specifying the moves, the induction on the ladder,
and the identification of that band.

By Lemma~\ref{lem:isotopy} and \eqref{eq:nfilling},
\[
       M_b\cong N\bigl(D_0:(1,-(b-5))\bigr)
                    \quad\text{as exteriors}.
\]
This homeomorphism is used to transfer hyperbolicity. The second
filling \eqref{eq:kfilling} uses the original disk marking of
\(M_b\). Twisting around either circle can change the expression
of the other circle's longitude, so no common longitude basis is
required for the two fillings.

\subsection{Three finite hyperbolicity certificates}\label{sec:certificates}

The computation certifies the complete hyperbolic structures of
\(N,M_5,M_6\), where \(M_5,M_6\) use \(Q_b\), with its full disk
twist, in the literal braid \(\cl{Q_bP_b}\). Table~\ref{tab:cert}
records the sizes of the input diagrams and certified triangulations.
The two direct tests for \(M_5,M_6\) give explicit bridge numbers
covered by the theorem regardless of the threshold obtained from \(N\).

\begin{table}[htbp]
\centering
\begin{tabular}{lrrr}
\toprule
Exterior & Diagram crossings & Cusps & Tetrahedra\\
\midrule
\(N\)   & 32 & 3 & 35\\
\(M_5\) & 32 & 2 & 31\\
\(M_6\) & 43 & 2 & 34\\
\bottomrule
\end{tabular}
\caption{Literal braid diagrams and complete ideal triangulations
for the three hyperbolicity certificates. Crossings are counted
in the given projection.}
\label{tab:cert}
\end{table}

Each exact input binds a planar-diagram code and its components
to the braid, checking every crossing port and every closure
identification. The exterior of that diagram is constructed in
Spherogram 2.4.1 and SnapPy 3.3.2. Its decorated triangulation
signature is checked against the retained one with the component
marking, and the face pairings and gluing equations of the
reloaded triangulation are checked. SnapPy under SageMath 10.9
then applies a Krawczyk interval test at 100 bits to both the
fresh exterior and the retained triangulation. For the analytic
criterion, a Krawczyk inclusion certifies a solution to a
square subsystem of the shape equations in the interval box.
The logarithmic edge and cusp equations are checked in the
correct branches, and every shape has positive imaginary part.
Together with the subsystem reduction, the positivity, gluing and
completeness checks give a complete
hyperbolic structure on the triangulated exterior
\cite[Sections 2--4]{HIKMOT}. Exact inputs, software pins,
rational interval bounds and the replay procedure are supplied
in the computational supplement. The proof uses these three
finite certificates and the written diagram induction.

\subsection{Proof of the prime-family conclusion}

Since \(N\) is complete and finite-volume hyperbolic, Thurston's
hyperbolic Dehn filling theorem
\cite[Theorem 5.8.2]{Thurston} applies to the \(D_0\) cusp with
the other two cusps unfilled. The slopes \((1,-(b-5))\) escape
every finite set as \(b\) grows. Thus \(M_b\) is hyperbolic for
every sufficiently large \(b\), as well as for \(b=5,6\) by
the direct certificates. For any such \(b\), apply the theorem
to \(M_b\) on the \(C\) cusp. The slopes in \eqref{eq:kfilling}
escape every finite set as \(k\) grows, so all sufficiently
large \(K_{b,k}\) are hyperbolic. A hyperbolic knot is prime.
Indeed, a composite knot group is the amalgam of the nontrivial
summand groups over their meridian, and the two summand longitudes
lie outside that amalgamating subgroup by peripheral injectivity.
They commute with the common meridian and fail to commute with
one another by amalgam normal form. The centralizer of a nonidentity
element of a torsion-free discrete subgroup of
\(\operatorname{PSL}_2(\mathbb C)\) is abelian, as is seen by
putting the element in parabolic or loxodromic normal form.

For a fixed such \(b\), the lengths \(\ell_k\) of the second
filling slopes on an embedded horotorus tend to infinity, since
they are the norms of \(\mu_C-d(d+1)k\lambda_C\) in a flat
lattice with \(\lambda_C\ne0\). For sufficiently large \(k\)
the fillings are hyperbolic and \(\ell_k>2\pi\).
The volume inequality of Futer, Kalfagianni and Purcell
\cite[Theorem 1.1]{FKP} and strict volume decrease under proper
hyperbolic filling \cite[Theorem 6.5.6]{Thurston} give
\[
 \left(1-\left(\frac{2\pi}{\ell_k}\right)^2\right)^{3/2}\!
     \operatorname{vol}(M_b)
   \ \le\ \operatorname{vol}\bigl(\Ext(K_{b,k})\bigr)
   \ <\ \operatorname{vol}(M_b).
\]
The volumes converge to \(\operatorname{vol}(M_b)\), and by
Mostow--Prasad rigidity each hyperbolic exterior type has
one volume. A finite set of exterior types would have a
maximum volume strictly below \(\operatorname{vol}(M_b)\).
Thus there are infinitely many knot types for that \(b\).
Together with the algebraic bounds, this proves
Theorem~\ref{thm:main}.
